% ============================================================================
%  CAPÍTULO 1 — INTRODUCCIÓN (req 4.2). Estructura obligatoria.
%  No incluye resultados ni conclusiones (req 4.3).
% ============================================================================
\chapter{Introducción}
\label{cap:introduccion}
\section{Antecedentes y motivación}
\label{sec:antecedentes}

\subsection{Contexto}
\label{sec:contexto}

La red hospitalaria pública chilena se organiza a través del Sistema Nacional de
Servicios de Salud (SNSS), que agrupa cerca de 192 establecimientos distribuidos
en 29 Servicios de Salud territoriales. La clasificación administrativa vigente
ordena estos hospitales en tres niveles de complejidad (alta, mediana y baja)
atendiendo a criterios estructurales: dotación de camas, disponibilidad de
especialidades médico-quirúrgicas y servicios de apoyo diagnóstico~\citep{Cid2024}.
Se trata de \emph{grupos administrativos predefinidos}: categorías
institucionales fijadas por infraestructura y no por la actividad clínica que
cada establecimiento efectivamente resuelve. Esta lógica estructural no captura
la heterogeneidad funcional real de la red, dado que hospitales con la misma
etiqueta pueden resolver cargas de enfermedad radicalmente
distintas~\citep{Havranek2023, Cid2024}.

Para cuantificar esa carga de enfermedad, el sistema ha adoptado de forma
progresiva los Grupos Relacionados por el Diagnóstico (GRD), un estándar
internacional que clasifica los egresos hospitalarios en categorías
clínicamente coherentes y con consumo de recursos similar a partir del
diagnóstico principal, los diagnósticos secundarios, los procedimientos
realizados y las características demográficas del paciente, vinculando así la
actividad clínica con los costos asociados~\citep{Cid2024}. Hacia 2024, su
implementación alcanzaba a 72 hospitales públicos de alta y mediana
complejidad~\citep{Cid2024}. Una descripción técnica detallada del sistema y de
su implementación en Chile se presenta en la Sección~\ref{sec:grd}.

Pese a esta riqueza de información clínica, los registros GRD han sido
aprovechados de manera limitada, con un uso predominantemente administrativo y
financiero orientado al pago por actividad~\citep{Cid2024}. Su estructura
homogénea, cobertura nacional y nivel de detalle ofrecen, sin embargo, una base
adecuada para ir más allá de ese uso administrativo y caracterizar de manera
objetiva la configuración funcional de la red.

En este trabajo, el \textbf{perfil funcional} de un hospital corresponde al
patrón de actividad asistencial observado en sus registros GRD, representado
mediante indicadores de volumen, complejidad, severidad, composición
diagnóstica, procedimientos, modalidades de ingreso y egreso, y características
demográficas de los pacientes atendidos. Es importante precisar que
``funcional'' en este sentido no equivale a desempeño, eficiencia, calidad
asistencial, infraestructura o capacidad instalada: el perfil funcional describe
\emph{qué} atiende un hospital, no \emph{qué tan bien} lo hace ni con \emph{qué
recursos} cuenta para hacerlo.

\subsection{Motivación y brecha de investigación}
\label{sec:motivacion}

La ausencia de una caracterización basada en la actividad clínica real tiene
consecuencias concretas. Quedó expuesta tras la implementación del mecanismo de
pago GRD en 2020, cuando los grupos administrativos predefinidos utilizados
para asignar la tasa de pago presentaron alta variabilidad interna que impidió
un financiamiento equitativo. Como solución de contingencia se fijó en 2023 una
tasa base única para toda la red~\citep{Cid2024}, pero esa medida no
resuelve el problema de fondo: sigue sin existir una segmentación funcional de
los hospitales que distinga sus verdaderos perfiles de producción.

Desde una perspectiva metodológica, la evidencia nacional sobre segmentación
funcional hospitalaria es escasa y carece de un marco reproducible que integre
aprendizaje no supervisado con mecanismos robustos de limpieza, adaptados a un
ecosistema de datos de calidad heterogénea. La experiencia internacional en
sistemas complejos como el francés y el chino respalda la utilidad de la
segmentación funcional (o \emph{análisis de conglomerados}) sobre datos de
producción hospitalaria~\citep{Chrusciel2021, Liu2025}, pero su aplicación al
caso chileno permanece escasamente explorada, en particular en lo relativo a
controlar el ruido de codificación y severidad propio de los registros
administrativos.

Esta investigación surge de esa brecha: sin una caracterización funcional,
reproducible y estadísticamente rigurosa, no es posible determinar si las
diferencias observadas entre hospitales responden a ineficiencias de gestión o a
cargas de enfermedad estructuralmente distintas. Disponer de esa
caracterización permitiría comparar a los establecimientos sobre bases más
equivalentes y generar insumos analíticos que apoyen la comprensión de la
estructura funcional de la red asistencial pública.

\subsection{Contribución esperada}
\label{sec:impacto}

Esta investigación tiene como propósito generar una caracterización funcional de
los hospitales públicos chilenos de alta y mediana complejidad a partir de los
registros GRD del período 2019--2024. Su principal resultado es un conjunto de
perfiles asistenciales emergentes (grupos de establecimientos con casuística
similar, junto a la interpretación de los atributos que los definen) y un
listado de hospitales con comportamiento atípico respecto de esos perfiles.

Los resultados de este trabajo constituyen un insumo exploratorio para análisis
posteriores de gestión hospitalaria, no una herramienta validada para auditoría
o asignación de recursos. En particular, dado que el estudio no incorpora
variables de costos, infraestructura, dotación de personal ni resultados
clínicos, cualquier uso de estos perfiles para decisiones de financiamiento o
gestión operativa requeriría complementarlos con esa información adicional. Su
aporte principal es más modesto y más acotado: proveer evidencia estructurada y
reproducible sobre la configuración funcional de la red, basada en datos
hospitalarios estandarizados, un enfoque que actualmente no existe a nivel
nacional. La comparación de los perfiles obtenidos con la clasificación
administrativa vigente y con la evidencia nacional e internacional permite,
además, identificar convergencias y divergencias que pueden motivar futuras
líneas de investigación.

\section{Descripción del problema}
\label{sec:problema}

Como se estableció en la Sección~\ref{sec:antecedentes}, no existe hoy una
caracterización funcional de los hospitales públicos chilenos que permita
distinguir, con base en su producción clínica real, si la posición relativa de
un establecimiento responde a ineficiencias de gestión o a cargas de enfermedad
estructuralmente distintas. Esta sección describe las dos dificultades
concretas que deben resolverse para construir esa caracterización a partir de
los datos disponibles: la calidad heterogénea de la fuente y su escala técnica.

La primera dificultad es la calidad heterogénea de los registros. El sistema
GRD, pese a reunir información clínico-administrativa detallada de las
hospitalizaciones, fue concebido con fines contables y presenta variaciones de
codificación, estructura y trazabilidad entre años que pueden introducir sesgos
silenciosos en cualquier análisis comparativo. La segunda es un desafío técnico
de escala: la base de datos consolidada supera los cinco millones de registros y
contiene más de un centenar de variables por egreso, volumen que exige un
procesamiento sistemático e imposibilita la revisión manual.

El problema de investigación es, por tanto, la ausencia de una metodología
integrada, reproducible y estadísticamente rigurosa que transforme estos datos
administrativos en una caracterización funcional de los hospitales, capaz de
aislar el ruido de codificación y severidad, distinguir perfiles asistenciales
reales e identificar establecimientos atípicos, más allá de la clasificación
administrativa vigente.

\subsection{Pregunta de investigación}
\label{sec:pregunta}

\begin{quote}
\textbf{¿Qué perfiles asistenciales y funcionales emergen al agrupar los hospitales 
públicos de alta y mediana complejidad según sus datos de producción GRD del 
período 2019--2024, y qué grado de concordancia y compactación en bloques de
variables retenidos presenta esa segmentación respecto de la clasificación
administrativa vigente?}
\end{quote}

De esta interrogante principal se desprenden las siguientes preguntas
secundarias:

\begin{enumerate}
  \item ¿Cuáles son los atributos dominantes de casuística y complejidad que
        definen la identidad de cada grupo, permitiendo una interpretación
        clínica y operativa de los perfiles funcionales generados?
  \item ¿Qué establecimientos presentan un comportamiento divergente respecto de
        los perfiles consolidados y cuál es la relevancia de su detección para
        garantizar un análisis comparativo equitativo?
  \item ¿Qué grado de concordancia o discrepancia presentan los grupos funcionales 
        emergentes al ser contrastados con la clasificación administrativa oficial 
        del Ministerio de Salud?
\end{enumerate}

\section{Solución propuesta}
\label{sec:solucion}

Frente a este problema, la solución propuesta consiste en un pipeline
reproducible de aprendizaje no supervisado que, a partir de los datos GRD
agregados por hospital del período 2019--2024, construye indicadores
institucionales de producción, complejidad, severidad, diversidad diagnóstica y
casuística clínica, y aplica un algoritmo de agrupamiento (\textit{clustering}
jerárquico) sobre dichos indicadores para obtener una segmentación funcional de
la red. La solución se complementa con la validación estadística de los grupos
obtenidos, la detección de hospitales atípicos y el contraste con la
clasificación administrativa del Ministerio de Salud. El detalle metodológico
de este pipeline se desarrolla en el Capítulo~\ref{cap:metodologia}.

\subsection{Hipótesis de trabajo}
\label{sec:hipotesis}

Se plantea como hipótesis de trabajo que una segmentación funcional basada en
indicadores GRD de producción, complejidad y casuística permite identificar
perfiles funcionales con cohesión y estabilidad interna evaluables en los
hospitales públicos chilenos de alta y mediana complejidad. La estabilidad se
evalúa exclusivamente como propiedad de la solución Ward bajo submuestreo. Su
relación con la clasificación administrativa vigente se examina mediante medidas
de concordancia y de compactación evaluadas en variables no utilizadas para
construir cada partición, sin interpretar el contraste como una prueba de
superioridad entre clasificaciones.

Esta hipótesis se sustenta en cuatro antecedentes desarrollados en el
Capítulo~\ref{cap:marco}: en primer lugar, la evidencia nacional que reporta baja coherencia
interna de la clasificación administrativa del MINSAL respecto de la
casuística real de los hospitales~\citep{Carlier2020, Giglio2021,
Villalobos2016}; en segundo lugar, la evidencia internacional de tipologías hospitalarias
construidas a partir de datos de producción~\citep{Chrusciel2021, Liu2025}; en tercer lugar, la
disponibilidad, en los registros GRD chilenos, de variables capaces de
representar casuística y complejidad clínica de forma estandarizada, y finalmente, la
escasez de estudios que integren, para el caso chileno, dimensiones clínicas
junto con criterios explícitos de estabilidad y detección de atipicidad. Los
criterios operacionales concretos para evaluar homogeneidad y estabilidad, así
como las pruebas estadísticas empleadas para contrastar la hipótesis, se
definen en el Capítulo~\ref{cap:metodologia}.

\section{Objetivos y alcances del proyecto}
\label{sec:objetivos}

\subsection{Objetivo general}
\label{sec:objetivo-general}

Caracterizar los perfiles funcionales de los hospitales públicos chilenos de
alta y mediana complejidad, con fines de comparabilidad institucional e
identificación de casos atípicos, mediante aprendizaje no supervisado sobre
datos GRD del período 2019--2024.

\subsection{Objetivos específicos}
\label{sec:objetivos-especificos}

\begin{enumerate}
  \item Construir una matriz institucional comparable de producción,
        complejidad y casuística clínica, mediante la depuración, integración
        y agregación de los registros GRD del período 2019--2024.
  \item Determinar una segmentación funcional de los hospitales, mediante
        agrupamiento jerárquico y criterios predefinidos de cohesión,
        separación y estabilidad interna bajo submuestreo.
  \item Caracterizar los perfiles obtenidos en términos de sus diferencias
        clínicas y operativas, mediante indicadores descriptivos y análisis
        post-hoc.
  \item Identificar los establecimientos atípicos respecto de los perfiles
        funcionales, mediante criterios multivariados de singularidad.
  \item Contrastar la segmentación funcional con la clasificación
        administrativa del MINSAL, mediante medidas de concordancia y de
        compactación evaluada en bloques de variables retenidos, sin establecer
        superioridad entre clasificaciones.
\end{enumerate}

\subsection{Alcances}
\label{sec:alcances}

El estudio se circunscribe a los hospitales públicos chilenos de alta y mediana
complejidad con presencia suficiente en el período 2019--2024, entendida como al
menos tres de los seis años del período con registros GRD. Se excluyen los
establecimientos de baja complejidad, dado que no generan datos compatibles con
el estándar técnico GRD.

El análisis se realiza sobre datos GRD agregados a nivel institucional. Si bien la
fuente primaria posee detalle por paciente, los indicadores se construyen por
hospital, por lo que el trabajo no desarrolla modelos para casos clínicos
individuales ni contempla el seguimiento de pacientes. No se incorpora
información de costos individuales ni datos clínicos más allá de los disponibles
en la base GRD, y la dotación de camas se excluye deliberadamente, de modo que el
perfilamiento se construye exclusivamente con información derivada de la actividad
clínica registrada. Por tratarse de registros administrativos, los resultados
están condicionados por la calidad de la codificación original, de modo que
cualquier sesgo o inconsistencia sistemática en los datos de entrada podría
repercutir en la precisión de la caracterización final.

El alcance final comprende la entrega de los perfiles funcionales o grupos y la
identificación de establecimientos atípicos como insumo analítico para la
caracterización de la red hospitalaria. Queda explícitamente fuera del alcance
cualquier recomendación sobre mecanismos de financiamiento, asignación
presupuestaria o gestión operativa de los hospitales, dado que estas decisiones
requieren incorporar variables de costos, infraestructura y recursos humanos que
no fueron incluidas en el análisis. Los resultados buscan aportar evidencia
empírica sobre la estructura funcional de la red, sin prescribir intervenciones
concretas de política pública o administrativa.


\section{Metodologías y herramientas utilizadas}
\label{sec:metodologias}

Para abordar el problema se empleó un flujo de análisis basado en el marco
CRISP-DM (\textit{Cross-Industry Standard Process for Data Mining})
~\citep{Chapman2000}, que integra la comprensión de los datos, la
preparación y reducción de dimensionalidad, el modelado mediante
aprendizaje no supervisado, la validación estadística de los grupos y la
detección de anomalías. El trabajo se implementa íntegramente en Python,
con un pipeline documentado y reproducible. La justificación conceptual de
este marco frente a alternativas, el detalle de las herramientas de
desarrollo empleadas y los parámetros de reproducibilidad del pipeline se
presentan en la Sección~\ref{sec:crisp-dm} y en el
Capítulo~\ref{cap:metodologia}.

\section{Organización del documento}
\label{sec:organizacion}

El presente documento se organiza en cinco capítulos que abordan de manera
progresiva el desarrollo de la investigación. A continuación se describe
brevemente el contenido de cada uno.

\begin{itemize}

  \item \textbf{Capítulo 1. Introducción.}
        Se presenta el contexto y la motivación del estudio, destacando la
        desconexión entre las clasificaciones administrativas vigentes y la
        realidad funcional de los hospitales públicos chilenos. Se describe el
        problema de investigación, se formula la hipótesis de trabajo, se
        definen los objetivos y se delimita el alcance del estudio.

  \item \textbf{Capítulo 2. Antecedentes.}
        Se revisan los conceptos fundamentales del sistema GRD y su implementación
        en Chile, las clasificaciones clínicas CIE-10 y CIE-9-MC, la experiencia
        internacional en perfilamiento hospitalario mediante clustering, y los
        antecedentes nacionales del Departamento de Ingeniería Informática de la
        USACH. Se presentan también las técnicas computacionales empleadas:
        aprendizaje no supervisado, reducción de dimensionalidad, detección de
        anomalías, validación de soluciones e interpretabilidad de modelos.

  \item \textbf{Capítulo 3. Metodología.}
        Se describe el diseño metodológico siguiendo el marco CRISP-DM. Se
        presentan las fuentes de datos, los criterios de elegibilidad de
        hospitales, el proceso de extracción y construcción de variables a nivel
        institucional, el preprocesamiento y normalización, y el esquema de
        clustering jerárquico y validación estadística aplicado.

  \item \textbf{Capítulo 4. Resultados y análisis.}
        Se presentan y discuten los perfiles asistenciales emergentes identificados
        por el pipeline, incluyendo la caracterización clínica de cada grupo, la
        detección de establecimientos atípicos y el contraste con la clasificación
        administrativa del Ministerio de Salud.

  \item \textbf{Capítulo 5. Conclusiones.}
        Se exponen las conclusiones finales del estudio, evaluando el cumplimiento
        de los objetivos planteados y los principales aportes de la investigación.
        Se discuten las limitaciones del trabajo y se proponen líneas de
        investigación futura.

\end{itemize}
